\documentclass{lecturenotes}

\usepackage{lecture_notes}
\usepackage{tikz}
\usetikzlibrary{decorations.pathreplacing}
\usepackage{tikz-cd}
% Reset mathcal back to previous style (not curly)
\DeclareMathAlphabet{\mathcal}{OMS}{cmsy}{m}{n}



%%%%%%%%%%%%%%%%%%%%%%%%%%%%%%%%%%%%%%%%%%%%%%%%%%%%%%%%%%%%%%%%%%%%%%%%%%%%%%%%%%%%%%%%%%%%%%%%%%%
%										For leaving comments									  %
%%%%%%%%%%%%%%%%%%%%%%%%%%%%%%%%%%%%%%%%%%%%%%%%%%%%%%%%%%%%%%%%%%%%%%%%%%%%%%%%%%%%%%%%%%%%%%%%%%%

%Name:			XXX
%Description:	Adds a piece of text in blue, surrounded by [] brackets. 
%				#1: the person the comment is addressed to
%				#2: the person the comment is from
%				#3: the comment
%Usage:			Write \XXX{collaborator}{myName}{myComment} to write [myComment] addressed to [collaborator]
\newcommand{\XXX}[3]{{\color{blue} \textbf{ [#1:  #3 \textit{ -#2-} ]}}}

%%%%%%%%%%%%%%%%%%%%%%%%%%%%%%%%%%%%%%%%%%%%%%%%%%%%%%%%%%%%%%%%%%%%%%%%%%%%%%%%%%%%%%%%%%%%%%%%%%%
%									     BibTeX commands 									      %
%%%%%%%%%%%%%%%%%%%%%%%%%%%%%%%%%%%%%%%%%%%%%%%%%%%%%%%%%%%%%%%%%%%%%%%%%%%%%%%%%%%%%%%%%%%%%%%%%%%

%Name:			Swap
%Description:	Command for properly typesetting "van der" and related expressions in Dutch and German names in the
%				bibliography. It swaps [van der] and [Lastname] in [van der Lastname] so that ordering will be performed on 
%				[Lastname] instead of [van der].
%Usage:			Write \swap{Lastname}{~van~der~}, Firstname instead of [van der Lastname, Firstname] in the author header
%				of the bibtex entry.

\newcommand*{\swap}[2]{\hspace{-0.5ex}#2#1}


%%%%%%%%%%%%%%%%%%%%%%%%%%%%%%%%%%%%%%%%%%%%%%%%%%%%%%%%%%%%%%%%%%%%%%%%%%%%%%%%%%%%%%%%%%%%%%%%%%%
%											Document											  %
%%%%%%%%%%%%%%%%%%%%%%%%%%%%%%%%%%%%%%%%%%%%%%%%%%%%%%%%%%%%%%%%%%%%%%%%%%%%%%%%%%%%%%%%%%%%%%%%%%%


\begin{document}

\begin{center}
	{\huge Additional Assessment 2MBA70 \\}
	\vspace*{10pt}
	{\large Date ??}
\end{center}

\vspace*{10pt}

This assessment consists of 4 problems reflecting results stated in the lecture notes. For solving them you are of course not allowed to simply use the corresponding result from the lecture notes. You are allowed to use any other results from the course and lecture notes. 

\vspace*{20pt}

\textbf{Problem 1 [3.2]}

Let $(\Omega, \cF)$ and $(E, \cG)$ be two measurable spaces such that $\cG = \sigma(\cA)$. Let $f: \Omega \to E$ be a function such that $f^{-1}(A) \in \cF$ for all $A \in \cA$.

%For this we have to show that $f^{-1}(B) \in \cF$ for each $B \in \cG$. We will do this by proving the collection of set $B$ that satisfy this property is equal to $\cG$.

\begin{enumerate}[label=(\alph*)]
\item Consider the following collection of subsets:
\[
	\cH := \{B \subset \cG \, : \, f^{-1}(B) \in \cF\}.
\]
Show that $\cH$ is a \sigalg/ on $E$.
\item Prove that $f$ is $(\cF, \cG)$-measurable. [Hint: use the inclusion property of generated \sigalgs/ to show that $\cG \subseteq \cH$.]
\end{enumerate}

\textbf{Problem 2 [4.5]}

Let $(\Omega, \cF, \mu)$ be a measure space, $f, g : \Omega \to \bbR_+$ two non-negative, measurable functions and $\alpha \geq 0$ be a constant. 

\begin{enumerate}[label=(\alph*)]
\item Suppose $B \in \cF$ satisfies $\mu(B) = 0$. Prove that
		\[
		\int_{B} f\, \dd \mu = 0. 
		\]
\item Suppose $f \leq g$. Prove that
		\[
		\int_\Omega f\, \dd \mu \leq \int_\Omega g\, \dd \mu.
		\]
\item Prove that
		\[
			\alpha \int_\Omega f\, \dd \mu = \int_\Omega (\alpha f )\,\dd \mu.
		\]
\end{enumerate}

[General hint: use simple functions first.]


\textbf{Problem 3 [8.3]}

Let $(X_n)_{n \ge 1}$ and $X$ be random variables defined on the same probability space. 

Prove that $X_n \stackrel{\bbP}{\rightarrow} X$ if and only if
\[
	\lim_{n \to \infty} \bbP(\|X_n - X\| > \varepsilon) = 0 \quad \text{for every } \varepsilon > 0.
\]

\pagebreak

\textbf{Problem 4 [11.2]}

Let $(\Omega, \cF, \bbP)$ be a probability space and $\cH$ be a sub-\sigalg/ of $\cF$. Furthermore, let $X, Y$ be random variables and $a \in \bbR$. 

\begin{enumerate}[label={(\alph*)}]
\item If $X$ is $\cH$-measurable,v prove that $\bbE[X | \cH] = X$ holds $\bbP$-almost everywhere.
\item Prove that $\bbE[a X | \cH] = a \bbE[X | \cH]$ holds $\bbP$-almost everywhere.
\item $\bbE[X + Y | \cH] = \bbE[X | \cH] + \bbE[Y | \cH]$ holds $\bbP$-almost everywhere.
\item If $X \le Y$, prove that $\bbE[X | \cH] \le \bbE[Y | \cH]$ holds $\bbP$-almost everywhere.
\end{enumerate}

\end{document}